\newpage
\chapter{Linear Systems of Equations}

\section{Linear Systems Standard Form}

A linear system of equations is a set of equations that can be expressed in the form:

\begin{align*}
    a_{1,1}x_1 + a_{1,2}x_2 + \cdots + a_{1,n}x_n & = b_1  \\
    a_{2,1}x_1 + a_{2,2}x_2 + \cdots + a_{2,n}x_n & = b_2  \\
                                               & \vdots \\
    a_{m,1}x_1 + a_{m,2}x_2 + \cdots + a_{m,n}x_n & = b_m,
\end{align*}

where \( a_{i,j} \) are the coefficients of the variables \( x_j \), and \( b_i \) are the constants
on the right-hand side of the equations. The goal is to find the values of \( x_1, x_2, \ldots, x_n \)
that satisfy all equations simultaneously.

More formally, given \(M, N\) with \(f: M \to N\). For each \(n \in N\) is the corresponding 
\emph{linear system of equation} given by 

\[
    f(m) = n
\]

With the solution set 

\[
    f^{-1}(\{n\}) = \{m \in M \mid f(m) = n\}
\]

Wit the following types of solutions depending on the nature of \(f\)

\begin{itemize}
 
    \item \emph{At least one solution} if \(f\) is surjective. 
   
    \item \emph{One solution and no more} if \(f\) is injective/surjective. 

\end{itemize}

\section{Linear System as Matrices}

A linear system can be represented in matrix form as:

\[
    A x = b
\]

Where \(A\) is the coefficient matrix, \(x\) is the vector of variables, and \(b\) is the vector
of constants. The coefficient matrix \(A\) is an \( m \times n \) matrix, where \( m \) is the number
of equations and \(n\) is the number of variables. The vector \(x\) is an \( n \times 1 \) column
vector, and the vector \(b\) is an \( m \times 1 \) column vector.

\subsection{Solution Set of a Linear System}

In the context of Definition 6.1, the solution set \( L(A, b) \) of the system of equations associated
with \( (A, b) \) is defined as

\[
    L(A, b) := \{ x \in \field^n \mid Ax = b \}.
\]

\subsubsection{Reduced Echelon Form}

For a matrix \(A \in \field^{m \times n}\) with rows \(Z_1, \dots, Z_n\). 

We define \(i \in \underbar{m}\) the \(i\)-t \emph{echelon index} \(Ei_i(A)\) as 

\[
    Ei := \min \{j \in \underbar{n} \mid A_{i,j} \ne 0\}
\]

We define the \emph{echelon form} as the sequence 

\[
    Ei(A) := (Ei_1 (A), Ei_2 (A), \dots,Ei_s (A))
\]

If the index of all \(Ei_i(A)\) coincides with the column number, then we call it the \emph{strict echelon form}.

\subsection{Equality of the Rank \texorpdfstring{\(\rank(A) = \rank_S(A)\)}{}}

We have \(\rank(A) = \rank_S(A)\).

\textbf{Proof:}

For \(x = {(x_1, \ldots, x_n)}^T \in \field^n\), it holds that

\[
    Ax = \sum_{i=1}^{n} x_i a_i.
\]

It follows that

\[
    \rank(A) = \rank(L_A) = \dim(\img(L_A))
    = \dim\left(\{L_A(x) \mid x \in \field^n\}\right)
    = \dim\left(\{Ax \mid x \in \field^n\}\right)
\]

\[
    = \dim\left\{ \sum_{i=1}^{n} x_i a_i \,\middle|\, x_i \in K \right\}
    = \dim(L(a_1, \ldots, a_n)) = \rank_S(A).
\]

\subsection{Solvability of the Linear System \texorpdfstring{\(Ax = b\)}{}}

The linear system \(Ax = b\) is solvable if and only if

\[
    \rank(a_1, \ldots, a_n) = \rank(a_1, \ldots, a_n, b).
\]

More concisely, we write

\[
    \rank(A) = \rank(A, b).
\]

\textbf{Proof:}

We have

\[
    Ax = (a_1, \ldots, a_n)x = x_1 a_1 + \cdots + x_n a_n = b.
\]

Therefore, a solution vector \(x\) exists if and only if

\[
    b \in L(a_1, \ldots, a_n) \quad \iff \quad \rank(A) = \rank(A, b),
\]

where \(L(a_1, \ldots, a_n)\) denotes the linear span of \(a_1, \ldots, a_n\).

\subsection{Solution Set of the Linear System \texorpdfstring{\(Ax = b\)}{}}

Let \(x_s \in \field^n\) be a solution of \(Ax = b\). Then the solution set is given by

\[
    L(A, b) = x_s + \ker(A) = \{x_s + x \mid x \in \ker(A)\}.
\]

\textbf{Proof:}

\[
    x \in \ker(A) \;\iff\; Ax = 0 \;\iff\; A(x_s + x) = Ax_s + Ax = Ax_s = b.
\]

\subsection{Free Variables}

In a system of linear equations, a \emph{free variable} is a variable whose column in the coefficient
matrix does not correspond to a pivot column after row reduction. Free variables are not directly solved by
any equation in the system; instead, they can take any real value, and the values of the leading
(or dependent) variables are expressed in terms of them.

Free variables arise naturally in systems with infinitely many solutions and play a key role in
describing the solution set. Given a system of with \(p\) equations and \(q\) unknowns then if
\(q > p\) then there are \(q - p\) free variables.

\textbf{Example:}

Consider the system

\[
    \begin{cases}
        x_1 + 2x_2 - x_3 = 4 \\
        3x_1 + 6x_2 - 3x_3 = 12
    \end{cases}
\]

which row reduces to

\[
    \begin{bmatrix}
        1 & 2 & -1 \\
        0 & 0 & 0
    \end{bmatrix}.
\]

Here, the pivot is in the first column corresponding to \( x_1 \), while \( x_2 \) and \( x_3 \) are free
variables. Assigning parameters to \( x_2 \) and \( x_3 \), the general solution can be expressed in terms
of these free variables.

\subsection{Gaussian Elimination}

\emph{Gaussian elimination} is a method for solving linear systems by transforming the system into an upper
triangular form. The steps involved in Gaussian elimination are:

\begin{enumerate}

    \item \emph{Forward elimination}: Transform the system into an upper triangular form by eliminating the
          variables from the equations.

    \item \emph{Back substitution}: Solve for the variables starting from the last equation and substituting
          back into the previous equations.

\end{enumerate}

The forward elimination process involves performing row operations on the augmented matrix
\([A | b]\) to create zeros below the diagonal. The row operations include:

\begin{itemize}

    \item Swapping two rows

    \item Multiplying a row by a non-zero scalar

    \item Adding or subtracting a multiple of one row from another row

\end{itemize}

These steps are defined as linear maps applied to a system like \(f(m) = n\) to \((g \circ f)(m) = g(n)\).
We need to be able to find the respective inverses to get to the result \((g \circ f)^{-1}({g(n)}) = n\). This 
is accomplished via the \emph{elementary matrices}.

Once the matrix is in upper triangular form, back substitution is used to find the values of the
variables. The last equation gives the value of the last variable, which can then be substituted into
the previous equations to find the other variables.

\subsubsection*{Procedure}

Starting with \(Ax = b\)

\[  
    \left[
    \begin{array}{cccc|c}
        a_{1,1} &  a_{1,2} & \cdots & a_{1,n}  & b_1    \\
        a_{2,1} &  a_{2,2} & \cdots & a_{2,n}  & b_2    \\
        \vdots  &   \vdots & \ddots & \vdots   & \vdots \\
        a_{m,1} &  a_{m,2} & \cdots & a_{m,n}  & b_m
    \end{array}
    \right]
\]

For eliminating the elements underneath a pivot element we apply 

\[
    \text{Row}_n - \frac{a_{n, i}}{a_{n - 1, i}} \text{Row}_{n - 1},
\]

where \(n\) the index of the target row and \(i\) the index of the column.

We denote the coefficient \(\frac{a_{n, i}}{a_{n - 1, i}} = l_{n, i}\) 

This results into

\[  
    \left[
    \begin{array}{cccc|c}
        a_{1,1} &  a_{1,2} & \cdots & a_{1,n}  & b_1    \\
        0 &  a_{2,2} - a_{1,2}\frac{a_{2,1}}{a_{1,1}} & \cdots & a_{2,n} - a_{1,n}\frac{a_{2,1}}{a_{1,1}}  & b_2 - b_1 \frac{a_{2,1}}{a_{1,1}}   \\
        \vdots  &   \vdots & \ddots & \vdots   & \vdots \\
        0 & 0  & \cdots & \lambda &  \gamma
    \end{array}
    \right]
\]

With \(\lambda\) and \(\gamma\) being a complex sums of results of the other manipulations.

\subsection{Invariability of the Solution Set}

Let \( f : M \to N \) and \( n \in N \). For every bijective map 
\( g : N \to N \), the following holds for the fiber of \( f \) over \( n \):

\[
    f^{-1}(\{n\}) = (g \circ f)^{-1}(\{g(n)\}).
\]

\textbf{Proof:}

The first inclusion does not use the bijectivity of \( g \):

Let \( m \in f^{-1}(\{n\}) \), i.e.\ \( f(m) = n \). Then
\( g(f(m)) = g(n) \), i.e.\ \( m \in (g \circ f)^{-1}(\{g(n)\}) \).
For the reverse implication one needs a left inverse of \( g \),
which of course exists because \( g \) is bijective.

\QED 

\subsection{Gauss-Jordan Elimination}

Gauss-Jordan elimination is an extension of Gaussian elimination that transforms the matrix into
reduced row echelon form (RREF). In RREF, each leading entry in a row is 1, and all entries above and
below the leading entry are zeros. The steps involved in Gauss-Jordan elimination are:

\begin{enumerate}

    \item Forward elimination: Transform the system into an upper triangular form.

    \item Back substitution: Transform the upper triangular matrix into RREF by eliminating the entries
          above the leading 1s.

    \item Solve for the variables directly from the RREF matrix.

\end{enumerate}

The Gauss-Jordan elimination method is particularly useful for finding the inverse of a matrix, as it 7
can be applied to the augmented matrix \([A | I]\), where \(I\) is the identity matrix. If the left side
of the augmented matrix becomes \(I\), then the right side will be the inverse of \(A\).

\textbf{Example:}

Solve the following system of equations using Gaussian elimination:

\begin{align*}
    2x + 3y + z  & = 1 \\
    4x + y - z   & = 2 \\
    -2x + y + 3z & = 3
\end{align*}

\textbf{Solution:} The augmented matrix for the system is:

\[
    \begin{bmatrix}
        2  & 3 & 1  & | & 1 \\
        4  & 1 & -1 & | & 2 \\
        -2 & 1 & 3  & | & 3
    \end{bmatrix}
\]

Performing row operations to eliminate the variables, we can transform the matrix into upper triangular form:

\[
    \begin{bmatrix}
        1 & \frac{3}{2} & \frac{1}{2} & | & \frac{1}{2} \\
        0 & -5          & -3          & | & 0           \\
        0 & 0           & 1           & | & 1
    \end{bmatrix}
\]

Now, we can perform back substitution to find the values of \(x\), \(y\), and \(z\):

\begin{align*}
    z           & = 1                           \\
    -5y - 3z    & = 0 \implies y = -\frac{3}{5} \\
    2x + 3y + z & = 1 \implies x = \frac{1}{5}
\end{align*}

Thus, the solution to the system is:

\begin{align*}
    x & = \frac{1}{5}  \\
    y & = -\frac{3}{5} \\
    z & = 1
\end{align*}

\subsection{Homogeneous Linear Equations}

A linear equation is said to be \emph{homogeneous} if its constant term is zero. That is, it
can be written in the form:

\[
    a_1 x_1 + a_2 x_2 + \cdots + a_n x_n = 0
\]

Such equations always have at least the trivial solution \(x_1 = x_2 = \cdots = x_n = 0\).

\subsection{Particular Solution}

A \emph{particular solution} to a linear system of equations is a specific solution that satisfies the system.
It can be found using various methods, including substitution, elimination, or matrix methods. A
particular solution is not unique; there may be multiple particular solutions depending on the system.
A particular solution can be found by substituting specific values for the variables and solving for the
remaining variables. For example, in the system

\begin{align*}
    2x + 3y & = 5 \\
    4x - y  & = 1
\end{align*}

we can substitute \(x = 1\) into the first equation to find \(y\):

\begin{align*}
    2(1) + 3y & = 5 \\
    3y        & = 3 \\
    y         & = 1
\end{align*}

Thus, \((x, y) = (1, 1)\) is a particular solution to the system. However, this is not the only solution;
other values of \(x\) may yield different values of \(y\).

\subsection{General = Particular + Homogeneous}

The \emph{general solution} of a linear system of equations is the complete set of solutions that satisfy the
system. It can be expressed as the sum of a particular solution and the general solution of the associated
homogeneous system.

The general solution can be written as:

\[
    x = x_p + x_h
\]

Where \( x_p \) is a particular solution to the non-homogeneous system, and \( x_h \) is the general
solution to the homogeneous system. The homogeneous system is obtained by setting the right-hand side
of the equations to zero:

\[
    A x = 0
\]

The theorem says:

Any linear system's solution set has the form:

\[
    \left\{ \vec{p} + c_1\vec{\beta}_1 + \cdots + c_k\vec{\beta}_k \;\middle|\; c_1, \ldots, c_k \in
    \R \right\},
\]

where \(\vec{p}\) is a particular solution to the system, and the vectors \(\vec{\beta}_1,
\ldots, \vec{\beta}_k\) form a basis of the solution space to the corresponding homogeneous system.
The number \(k\) equals the number of \emph{free variables} the system has after applying Gaussian
elimination.

\textbf{Proof:}

We want to proof that every linear system of equations \(Ax = b\) which has at least one solution fulfills 
the following properties 

\begin{enumerate}

    \item \(L = \{Ax = 0\} \ne \emptyset\)

    \item The general solution is linear combination of the particular solutions plus the homogeneous solution.

\end{enumerate}

For the first property the solution of \(Ax = 0\) is the \(\ker(A)\) more precisely \(\ker(\phi_A)\) and thus, 
it is not empty because \(0 \in \ker(A)\)

The second proposition follows for the second one because the addition of the zero vector to a linear 
combination of other vectors just gives you the original combination by the axioms of the vector spaces.

\QED 

\subsection{Linear Combination Lemma}

Any linear combination of linear combinations is a linear combination.

\subsection{Example: Gaussian Elimination with 3 Equations and 4 Unknowns}

Consider the following system of linear equations:

\begin{align*}
    x_1 + 2x_2 + x_3 + x_4   & = 4  \\
    2x_1 + 5x_2 + x_3 + 3x_4 & = 10 \\
    x_1 + 3x_2 + 2x_3 + 2x_4 & = 7
\end{align*}

\textbf{Step 1: Augmented Matrix}

\[
    \begin{bmatrix}
        1 & 2 & 1 & 1 & 4  \\
        2 & 5 & 1 & 3 & 10 \\
        1 & 3 & 2 & 2 & 7  \\
    \end{bmatrix}
\]

\textbf{Step 2: Eliminate below pivot in column 1}

\begin{itemize}
    \item Row 2 \(=\) Row 2 \(-\) 2 \(\times\) Row 1
    \item Row 3 \(=\) Row 3 \(-\) Row 1
\end{itemize}

\[
    \begin{bmatrix}
        1 & 2 & 1  & 1 & 4 \\
        0 & 1 & -1 & 1 & 2 \\
        0 & 1 & 1  & 1 & 3 \\
    \end{bmatrix}
\]

\textbf{Step 3: Eliminate below pivot in column 2}

\begin{itemize}
    \item Row 3 \(=\) Row 3 \(-\) Row 2
\end{itemize}

\[
    \begin{bmatrix}
        1 & 2 & 1  & 1 & 4 \\
        0 & 1 & -1 & 1 & 2 \\
        0 & 0 & 2  & 0 & 1 \\
    \end{bmatrix}
\]

\textbf{Step 4: Back Substitution}

From Row 3:
\[
    2x_3 = 1 \implies x_3 = \frac{1}{2}
\]

From Row 2:
\[
    x_2 - x_3 + x_4 = 2 \implies x_2 = 2 + x_3 - x_4 = 2 + \frac{1}{2} - x_4 = \frac{5}{2} - x_4
\]

From Row 1:
\[
    x_1 + 2x_2 + x_3 + x_4 = 4
    \implies x_1 = 4 - 2x_2 - x_3 - x_4
\]
Substitute:
\[
    x_1 = 4 - 2\left( \frac{5}{2} - x_4 \right) - \frac{1}{2} - x_4
    = 4 - 5 + 2x_4 - \frac{1}{2} - x_4
    = -1 - \frac{1}{2} + x_4 = -\frac{3}{2} + x_4
\]

\textbf{General Solution}

Let \(x_4 = t\) (free variable), then:

\[
    \begin{aligned}
        x_1 & = -\frac{3}{2} + t \\
        x_2 & = \frac{5}{2} - t  \\
        x_3 & = \frac{1}{2}      \\
        x_4 & = t
    \end{aligned}
    \quad \text{with } t \in \R
\]

\textbf{Solution Set:}

\[
    \left\{
    \begin{bmatrix}
        -\frac{3}{2} \\ \frac{5}{2} \\ \frac{1}{2} \\ 0
    \end{bmatrix}
    + t \cdot
    \begin{bmatrix}
        1 \\ -1 \\ 0 \\ 1
    \end{bmatrix}
    \;\middle|\; t \in \R
    \right\}
\]

\subsection{The Determinant, the cross product and the solutions of linear Systems of Equations}

A linear system of three equation has the following properties:

\begin{itemize}

    \item There is a unique solution if the determinant of the coefficient matrix is non-zero.

          \[
              \langle (a \times b), c\rangle = \det(a,b,c) \ne 0
          \]

    \item There are infinitely many solutions if the determinant of the coefficient matrix is zero.

          \[
              \langle (a \times b), c\rangle = \det(a,b,c) = 0
          \]

    \item There is no solution if the determinant of the coefficient matrix is zero and the system is inconsistent.

          \[
              \langle (a \times b), c\rangle = 0
          \]

\end{itemize}

\subsection{Chart for the number of solutions}

\begin{tikzpicture}[node distance=1cm and 2cm]
    \node [terminator] (start) {Linear \(n \times n\) System};
    \node [decision, below=of start] (decision1) {\(\det(A) = 0?\)};
    \node [decision, below=of decision1] (decision2) {\(\rank(A) = \rank(A,b)?\)};
    \node [terminator, left=of decision1] (oneS) {Exactly one solution};
    \node [terminator, below left=of decision2] (inS) {Inf. sol. with \((n - \rank(A))\) parameters};
    \node [terminator, below=of decision2] (noS) {No solution};

    \path [connector] (start) -- (decision1);
    \path [connector] (decision1.west) -- node[above] {No} (oneS.east);
    \path [connector] (decision1) -- node[right] {Yes} (decision2);
    \path [connector] (decision2.south west) -- node[above left] {Yes} (inS.north east);
    \path [connector] (decision2.south) -- node[above right] {No} (noS.north);
\end{tikzpicture}

\bigskip

\begin{tikzpicture}[node distance=1cm and 2cm]
    \node [terminator] (start) {Linear \(m \times n\) System};
    \node [decision, below=of start] (decision1) {\(\rank(A) = \rank(A,b)?\)};
    \node [terminator, left=of decision1] (noS) {No solution};
    \node [decision, below=of decision1] (decision2) {\(\rank(A) = n?\)};
    \node [terminator, below left=of decision2] (inS) {Inf. sol. with \((n - \rank(A))\) parameters};
    \node [terminator, below right=of decision2] (oneS) {Exactly one solution};

    \path [connector] (start) -- (decision1);
    \path [connector] (decision1.west) -- node[above] {No} (noS.east);
    \path [connector] (decision1) -- node[right] {Yes} (decision2);
    \path [connector] (decision2.south west) -- node[above left] {No} (inS.north east);
    \path [connector] (decision2.south east) -- node[above right] {Yes} (oneS.north west);
\end{tikzpicture}

\subsection{Proof of the Gaussian Elimination}

The row operations do not change the solution set.

\textbf{Proof:}

A linear system of equations of the form \(Ax = b\), \(A \in \field^{m \times n}\) with some solution
\(x\) can be multiplied by some elementary matrix \(C_1, C_2, C_3\) which correspond to a row
operation \(CAx = Cb\).

The Equality is kept and therefore, \(L(A,b) \subset L(CA, Cb)\). If \(x \in L(CA, Cb)\) then

\[
    CAx = Cb \implies C^{-1}CAx = C^{-1}Cb \implies Ax = b
\]

Because of the invertibility of elementary matrices we can be sure that  \(L(CA, Cb) \subset L(A,b) \).

Now for the case that \(L(A,b) = \emptyset\). If we assume that \(L(CA,CB) \ne \emptyset \), then
there would be some \(x \in L(CA,CB)\) with \(CAx = Cb\). By multiplying both side by \(C^{-1}\) we
get \(Ax = b\) and therefore, \(L(A, b) \ne \emptyset \),  which is a contradiction to our original
assumption  \(L(A,b) = \emptyset\).

\QED

Now we only have to prove that each regular matrix can be transformed in the reduced echelon form or
even row reduced echelon form. This is the equivalent of saying that the row operations do not change
the rank of the matrix.

\textbf{Proof:}

If \(C\) is an elementary matrix and \(A \in \field^{m \times n}\) some matrix. The statement \(rg(CA) = rg(A)\)
is true because of the invertibility of the elementary matrices and the isomorphism in vector spaces.

Given \(A = (a_1, a_2, \dots, a_n)\). The exchange of two columns preserves \(L := L(a_1, \dots, a_n)\),
and also the \(rg_s (A)\). Concerning the scaling of a column by some \(\lambda \) \(L := L(a_1, \dots, \lambda a_k, \dots, a_n)\)
also preserves \(dim(L) = rg_s(A)\). And finally, the addition \(L = L(a_1, \dots, a_i + \lambda a_k, \dots, a_n)\) also preserves
\(rg_s (A) = rg(A)\). Thus, \(rg(CA) = rg(A)\).

\QED

The only part left is to prove that if \(A\) has an inverse, then with use of the row operations we can
transform it in to a matrix with an upper triangular with also no zero element on the main diagonal.

\textbf{Proof:}

First we eliminate all the elements under the main diagonal for all columns until the column \(k - 1\).
With this, we have created a new matrix \(A^(\ell) = C_{\ell} \dot \cdots \dot C_1 A \) which looks
like:

\[
    A^{(\ell)} =
    \begin{bmatrix}
        a^{(\ell)}_{1,1} & \cdots & a^{(\ell)}_{1,k-1}   & a^{(\ell)}_{1,k}    & a^{(\ell)}_{1,k+1}   & \cdots & a^{(\ell)}_{1,n}    \\
        \vdots          & \ddots & \vdots               & \vdots             & \vdots               & \ddots & \vdots             \\
        0               & \cdots & a^{(\ell)}_{k-1,k-1} & a^{(\ell)}_{k-1,k} & a^{(\ell)}_{k-1,k+1} & \cdots & a^{(\ell)}_{k-1,n} \\
        0               & \cdots & 0                    & a^{(\ell)}_{k,k}    & a^{(\ell)}_{k,k+1}   & \cdots & a^{(\ell)}_{k,n}    \\
        0               & \cdots & 0                    & \vdots             & a^{(\ell)}_{k+1,k+1} & \cdots & a^{(\ell)}_{k+1,n} \\
        \vdots          & \ddots & \vdots               & \vdots             & \vdots               & \ddots & \vdots             \\
        0               & \cdots & 0                    & a^{(\ell)}_{n,k}    & \cdots               & \cdots & a^{(\ell)}_{n,n}
    \end{bmatrix}
\]

Because of \(rg(A) = rg(A^{(\ell)})\), \(a^{(\ell)}_{i,i} \ne 0\) for \(1 \le i < k\). All of these
elements will not change for the rest of the Gaussian Elimination. Now we want to build another matrix
with row operations \(A^{(\mu)}\) where \(a^{(\mu)}_{k,k} \ne 0\) and \(a^{(\mu)}_{i,k} = 0\) for
\(i > k\). If \(a^{(\ell)}_{k,k} \ne 0\) then \((i,k)\)-th element will become 0 by subtracting
\(\frac{a^{(\ell)}_{ink}}{a^{(\ell)}_{k,k}}\) time the \(k\)-th row of the \(i\)-th row. So can the \(k\)-th
column be transformed with a series of row operations to the desired form. But if \(a^{(\ell)}_{k,k} = 0\),
then you have to swap it with another row under it. This is the critical point of the algorithm: What
if there is no row, also \(a^{(\ell)}_{i,k} = 0\) for \(i \ge k\)? Then \(A^{(\ell)}\) would be:

\[
    A^{(\ell)} =
    \begin{bmatrix}
        a^{(\ell)}_{1,1} & \cdots & a^{(\ell)}_{1,k-1}   & a^{(\ell)}_{1,k}    & a^{(\ell)}_{1,k+1}   & \cdots & a^{(\ell)}_{1,n}    \\
        \vdots          & \ddots & \vdots               & \vdots             & \vdots               & \ddots & \vdots             \\
        0               & \cdots & a^{(\ell)}_{k-1,k-1} & a^{(\ell)}_{k-1,k} & a^{(\ell)}_{k-1,k+1} & \cdots & a^{(\ell)}_{k-1,n} \\
        0               & \cdots & 0                    & 0                  & a^{(\ell)}_{k,k+1}   & \cdots & a^{(\ell)}_{k,n}    \\
        0               & \cdots & 0                    & \vdots             & a^{(\ell)}_{k+1,k+1} & \cdots & a^{(\ell)}_{k+1,n} \\
        \vdots          & \ddots & \vdots               & \vdots             & \vdots               & \ddots & \vdots             \\
        0               & \cdots & 0                    & 0                  & \cdots               & \cdots & a^{(\ell)}_{n,n}
    \end{bmatrix}
\]

But, that would imply that the \(k\)-th column is a linear combination of the other \(k - 1\) columns,
and then \(rg(A^{(\ell)}) \le n - 1\) which is a contradiction to the assumption  \(rg(A^{(\ell)}) = n\)

\QED

\subsection{Alternative Method for finding the Rank of a matrix}

For a matrix \(A \in \field^{(m \times n)}\). The biggest submatrix of the size \(r \times r\) with a
non-zero determinant gives us the \emph{rank} of the matrix. \(rg(A) = r\).

\subsection{Cramer's Rule}

For a square matrix \(A = (a_1, \dots, a_n)\) and \(x,b \in \field^{n}\) with \(Ax = b\) with \(\det A \ne 0\)
then

\[
    A_i = (a_1, \dots, a_{i - 1}, b, a_{i + 1}, \dots, a_n), i = 1,\dots,n.
\]

The solution \(i\) is then

\[
    x_i = \frac{\det(A_i)}{\det(A)}
\]

\textbf{Proof:}

Given a matrix \(X_i \in \field^{n \times n}\) by

\[
    \begin{bmatrix}
        1      &   & \cdots    & x_1    &  & \cdots \\
        \vdots &   & \ddots    & \vdots &  & 0      \\
               &   & x_{i - 1} &        &  &        \\
               & 0 & x_i       &        &  &        \\
               &   & x_{i + 1} &        &  &        \\
               &   & \vdots    & \ddots &  &        \\
               &   & x_n       &        &  & 1
    \end{bmatrix}
\]

\(X_i = (e_1, \dots, e_{i - 1}, e_i, e_{i + 1}, \dots, e_n)\). By using an \(n\)-times the
Laplace formula in the first row shows \(|X_i| = x_i\).

\begin{align*}
    A X_i & = A (e_1, \dots, e_{i - 1}, e_i, e_{i + 1}, \dots, e_n)     \\
    A X_i & = A (a_1, \dots, a_{i - 1}, Ax, a_{i + 1}, \dots, a_n)      \\
    A X_i & = A (a_1, \dots, a_{i - 1}, b, a_{i + 1}, \dots, a_n) = A_i
\end{align*}

This implies \(|A_i| = |AX_i| = |A||X_i|\) therefore, \(x_i = \frac{|A_i|}{A}\).

\QED

Another way of thinking about this in a more geometric way is to first think at the determinant as
the area/volume/etc. that is spanned by the basis vector and the stretching factor of the transformation
\(\det(A)\). Now visualize your matrix \(A\) as a linear transformation with
a non-zero determinant and think about \(\det(A_i)\) as the area/volume/etc. which is spanned by the
basis vector but with one of the vector substituted by \(b\).

Now in two dimensions, think of the area spanned by the basis vectors \(A = xy\).
The signed area of the original parallelogram gets
stretched by \(\det(A)\) this gives us \(Area = \det(A)y \implies y = \frac{Area}{\det(A)}\). This
also works for \(x\).

\subsection{LU Decomposition}

Given a regular system \(Ax = b\) with \(a_{k,k}^{(k - 1)} \ne 0\), then the Gaussian elimination delivers a so called 
\emph{\(LU\)-Decomposition} 

\[
    A = LU
\]

of \(A\), where \(L\) is a under triangular matrix with main diagonal consisting of only of ones and 
\(U\) being a upper triangular matrix. \(a_{k,k}^{(k - 1)} \ne 0\) depicts the elements of the main diagonal for the 
\(k - 1\) matrix generated after applying a Gauss operation.

\subsubsection*{Derivation}

Given our matrix \(A\), we want to eliminate the coefficients underneath the main diagonal. This is accomplished by applying an elementary matrix of the form: 

\[
    F = 
    \begin{bmatrix}
        1 &        &         &   &        \\
          & \ddots &         &   &        \\
          &        & 1       &   &        \\
          &        & \lambda & 1 &        \\
          &        &         &   & \ddots \\
    \end{bmatrix}
\]

Hence, to get to the row-echelon-form we apply a series of such \(F\) matrices, where for the 
elimination of the \(k\)-th row we use the \emph{Frobenius} matrix 

\[
    F_k = 
    \begin{bmatrix}
        1 &        &         &   &        \\
          & \ddots &         &   &        \\
          &        & 1       &   &        \\
          &        & -l_{k +1, k} & 1 &        \\
          &        &  -l_{n,k}       &   & \ddots \\
    \end{bmatrix}
\]

where \(-l_{i,k} = \frac{a_{i,k}^{(k - 1)}}{a_{k,k}^{k - 1}}\). 

For example: 

\[
    A = 
    \begin{bmatrix}
        a_{1,1} & a_{1,2} \\ 
        a_{2,1} & a_{2,2}
    \end{bmatrix}
    ,
    a_{1,1} \ne 0
\]

To eliminate \(a_{2,1}\) we need to subtract \(\frac{a_{2,1}}{a_{1,1}}\) times the first row.

\[
    F = 
    \begin{bmatrix}
        1 & 0 \\ 
       - \frac{a_{2,1}}{a_{1,1}} & 1
    \end{bmatrix}
    = 
    \begin{bmatrix}
        1 & 0 \\ 
        - l_{2,1} & 1
    \end{bmatrix}
\]

The multiplication results into

\[
    FA = 
    \begin{bmatrix}
        a_{1,1} & a_{1,2} \\ 
        a_{2,1}  - \frac{a_{1,1}a_{2,1}}{a_{1,1}} & a_{2,2} - \frac{a_{2,1}}{a_{1,1}}
    \end{bmatrix}
    = 
    \begin{bmatrix}
        a_{1,1} & a_{1,2} \\ 
        0 & a_{2,2} - \frac{a_{2,1}}{a_{1,1}}
    \end{bmatrix}
\]

Note that the main diagonal does not have to be normalized.

By repeating this process we get 

\[
    U = A^{(n - 1)} = F_{n - 1} \cdot \cdots \cdot F_1 A
\]

and 

\[
    y = b^{(n - 1)} = F_{n - 1} \cdot \cdots \cdot F_1 b
\]

Due to all \(F_k\) being regular, we get that \(Ax = b\) and \(Ux = y\) are equivalent.
Hence, by applying the inverses 

\[
    A = F_{1}^{n - 1} \cdot \cdots \cdot F_{n - 1} U
\]

This chain of the inverses is called 

\[
    L = 
    \begin{bmatrix}
        1 &        &         &   &        \\
        l_{2,1}  & \ddots &         &   &        \\
        \vdots &        & 1       &   &        \\
          &        &  & 1 &        \\
        l_{n,1} &        &  \cdots      &  l_{n,n-1} & 1 \\
    \end{bmatrix}
\]

The individual members are 

\[
    F_{k}^{-1} = 
    \begin{bmatrix}
        1 &        &         &   &        \\
          & \ddots &         &   &        \\
          &        & 1       &   &        \\
          &        & l_{k+1,k} & 1 &        \\
          &        & l_{n,k}      &   & 1 \\
    \end{bmatrix}
\]

Therefore \(A = LU\) and \(Ly = b\).

\subsubsection{Algorithm}

Overview: 

\begin{enumerate}
    
    \item Get \(L\) and \(U\) via Gaussian Elimination. 

    \item Solve \(Ly = b\). 

    \item Solve \(Ux = y\). 

\end{enumerate}

Concrete algorithm: 

Start with the your \(n \times n\) matrix perform the following steps across all of the 
entries in the main diagonal and its corresponding sub-square-matrix.

\begin{enumerate}

    \item \textbf{Division}: For the next \(k+n\) rows divide the elements in the current column by \(a_{k,k}\). Only to apply to the \(k \times k\) matrix.

    \item \textbf{Column Elimination}: Subtract from the current row \(a_{k + n}\)-times the \(k\) row, also your current first row. Apply only for the inner sub-matrix \((k - 1) \times (k - 1)\)

    \item \textbf{Scoping}: Now proceed with the next diagonal element which is going to give you an inner \((n - 1) \times (n - 1)\) matrix. To repeat 
    the process. Note that when performing the row addition, this will only affect the inner matrix; hence, our old results remain intact.
    
    \item \textbf{Getting \(L\) and \(U\)}: After the previous process, the resulting matrix will be divided into two matrices. \(L\) consisting on 
    the elements underneath the main diagonal with it being all ones and above all zeros; and \(U\) begin the rest upper triangular matrix.
\end{enumerate}

\textbf{Example:}

Given

\[
    A = 
    \begin{bmatrix}
        2 & 1 & 1 \\ 
        4 & 4 & 0 \\ 
        2 & 7 & 1
    \end{bmatrix}, 
    b = 
    \begin{bmatrix}
        7 \\ 
        12 \\ 
        19
    \end{bmatrix}
    \text{ and }
    b' = 
    \begin{bmatrix}
        3 \\ 
        4 \\ 
        9
    \end{bmatrix}
\]

solve using the \(LU\)-decomposition.

First iteration:

\[
    \begin{bmatrix}
        2 & 1 & 1 \\ 
        4 & 4 & 0 \\ 
        2 & 7 & 1
    \end{bmatrix}  
    \rightarrow 
    \begin{bmatrix}
        2 & 1 & 1 \\ 
        2 & 4 & 0 \\ 
        1 & 7 & 1
    \end{bmatrix}
    \rightarrow 
    \begin{bmatrix}
        2 & 1 & 1 \\ 
        2 & 2 & -2 \\ 
        1 & 6 & 0
    \end{bmatrix}
\]

Second iteration; 

\[
    \begin{bmatrix}
        2 & 1 & 1 \\ 
        2 & 2 & -2 \\ 
        1 & 6 & 0
    \end{bmatrix}
    \rightarrow 
    \begin{bmatrix}
        2 & 1 & 1 \\ 
        2 & 2 & -2 \\ 
        1 & 3 & 0
    \end{bmatrix}
    \rightarrow
    \begin{bmatrix}
        2 & 1 & 1 \\ 
        2 & 2 & -2 \\ 
        1 & 3 & 6
    \end{bmatrix}
\]

Hence 

\[
    L = 
    \begin{bmatrix}
        1 & 0 & 0 \\ 
        2 & 1 & 0 \\ 
        1 & 3 & 1
    \end{bmatrix}
    \quad 
    U = 
    \begin{bmatrix}
        2 & 1 & 1 \\ 
        0 & 2 & -2 \\ 
        0 & 0 & 6
    \end{bmatrix}
\]

We solve \(Ly = b\)

\[
    \begin{bmatrix}
        1 & 0 & 0 \\ 
        2 & 1 & 0 \\ 
        1 & 3 & 1
    \end{bmatrix}
    y 
    = 
    \begin{bmatrix}
        7 \\ 
        12 \\ 
        19
    \end{bmatrix}
    \iff 
    y = 
    \begin{bmatrix}
        7 \\ 
        -2 \\ 
        18
    \end{bmatrix}
\]

Then we solve \(Ux = y\)

\[
    \begin{bmatrix}
        2 & 1 & 1 \\ 
        0 & 2 & -2 \\ 
        0 & 0 & 6
    \end{bmatrix}
    x 
    = 
    \begin{bmatrix}
        7 \\ 
        -2 \\ 
        18
    \end{bmatrix} 
    \iff
    x = 
    \begin{bmatrix}
        1 \\ 
        2 \\ 
        3
    \end{bmatrix} 
\]

\subsubsection{Extended Algorithm for Any Regular Matrix}

We want to expand the algorithm to be compatible with the swap rows operations. Given a regular matrix 
\(A^{(k - 1)} = F_{k - 1} \cdot \dots \cdot F_1 A\), where after a Gaussian operations we get \(a_{k,k}^{(k - 1)} = 0\), then 
there exists in the \(k\)-th. row underneath at a least a \(a_{i,k}^{(k - 1)} \ne 0, i > k\), so that we can swap the rows.

\textbf{Proof:}

Let us suppose, that there does not exists such an element; then the \(k-th\) column must be a linear combination of the other 
columns because there is a free variable, which contradicts our initial assumption. Note that due to \(F_{i}\) also being regular, the multiplication 
of \(A^{(k - 1)}\) must also be regular. Therefore, it is always possible to swap rows.

\QED

\subsubsection{Derivation of the Algorithm}

We start with the same setup as before, but now we also have 

\[
    P = 
    \begin{bmatrix}
        1 &        &         &   &            &   &   &        & \\
          & \ddots &         &   &            &   &   &        & \\
          &        & 0       &   &            & 1  &   &        & \\
          &        &         & 1 &            &  &   &        & \\
          &        & \vdots  &   &  \ddots    & \vdots &   &        & \\
          &        &  1      &   &            & 0 &   &        & \\
          &        &         &   &            &   & 1 &        & \\
          &        &         &   &            &   &   & \ddots & \\
          &        &         &   &            &   &   &        &  1\\
    \end{bmatrix}
\]

which the following properties \(P = P^T\) and \(PP = I\) which means \(P = P^{-1}\)

Our \(U\) and \(L\) becomes 

\[
    U = F_{n - 1} P_{n - 1} \cdot \dots \cdot F_1 P_1 A
\]

\[
    \hat{L} = (F_{n - 1} P_{n - 1} \cdot \dots \cdot F_1 P_{-1} )^{-1} 
\]

And \(U = \hat{L}^{-1} A\), the question is how to get \(\hat{L}\) to always be a lower triangular matrix.

For that we use the inverse properties of permutation matrices to multiply by the identity matrix 

\[
    I = P_1 P_2 \cdot \dots \cdot P_{n - 1} P_{n - 1} \cdot \dots \cdot P_2 P_1
\]

by associativity. 

We get 

\begin{align*} 
    U &= F_{n - 1} P_{n - 1} \cdot \dots \cdot F_1 P_1 A \\ 
      &= F_{n - 1} P_{n - 1} \cdot \dots \cdot F_1 P_1 \cdot (P_1 P_2 \cdot \dots \cdot P_{n - 1} P_{n - 1} \cdot \dots \cdot P_2 P_1) A\\ 
      &= (F_{n - 1} P_{n - 1} \cdot \dots \cdot F_1 P_1  \cdot P_1 P_2 \cdot \dots \cdot P_{n - 1}) P_{n - 1} \cdot \dots \cdot P_2 P_1 A\\ 
      &= (F_{n - 1} P_{n - 1} \cdot \dots \cdot F_1 P_1  \cdot P_1 P_2 \cdot \dots \cdot P_{n - 1}) P A\\ 
      & = L^{-1} P A
\end{align*}

Where 

\begin{align*} 
    L &= (F_{n - 1} P_{n - 1} \cdot \dots \cdot F_1 P_1  \cdot P_1 P_2 \cdot \dots \cdot P_{n - 1})^{-1} \\ 
      &=  P_{n - 2} \cdot \dots \cdot P_2 F_{1}^{-1} P_2 F_{2}^{-1} \cdot \dots \cdot P_{n - 1} F_{n - 1}^{-1}
\end{align*}

Now \(L\) is lower triangular matrix, because of the following inductive proof: 

Let \(L_1 = F_{1}^{-1}\), \(L_k = P_k L_{k -1} P_k F_{k}^{-1}\) for \(k = 2, \dots, n - 1\). This is \(L_{n  - 1} = L\)

We show that all \(L_k\) are lower triangular matrices.

\emph{Base Case}: For \(k = 1\)

\[
    L_1 = 
        \begin{bmatrix}
            1 &        &         &   &        \\
             l_{2,1}& \ddots &         &   &        \\
             \vdots&        & 1       &   &        \\
            &        &  & 1 &        \\
            l-{n,1} &        &        &   & \ddots \\
        \end{bmatrix}
\]

\emph{Inductive Step}:

For \(L_2 = P_2 L_1 P_2 F_{2}^{-1}\): 

In \(P_2 L_1\) are the second and \(j\) row changed

\[
    P_2 L_1= 
        \begin{bmatrix}
            1 &        &         &   &        \\
             l_{2 1}& \ddots &         &  1 &        \\
             l_{j,1} &        & 1       &   &        \\
            & \vdots       &  & 1 &        \\
            l-{n,1} &        &        &   & \ddots \\
        \end{bmatrix}
\]

Then in \(P_2 L_1 P_2\) the change get undone; and finally, when multiplied by 

\[
    F_{2}^{-1} =     
        \begin{bmatrix}
            1 &        &         &   &        \\
              & \ddots &         &   &        \\
              &  l_{3,2}     & 1       &   &        \\
              &   \vdots    &  & 1 &        \\
              &  l_{n,2}      &        &   & \ddots \\
        \end{bmatrix}
\]

we get 

\[
    L_2 = P_2 L_1 P_2 F_{2}^{-1} =     
        \begin{bmatrix}
            1 &        &         &   &        \\
            l_{2,1}  & \ddots &         &   &        \\
            \vdots  &  l_{3,2}     & 1       &   &        \\
              &   \vdots    &  & 1 &        \\
            l_{n, 1}  &  l_{n,2}      &        &   & \ddots \\
        \end{bmatrix}
\]

Hence, by induction \(L_k\) is always a lower triangular matrix.

\QED 

Formally we state that given a regular matrix \(A \in \R^{n \times n}\), there are also permutations \(P_k, k = 1, \dots n - 1\) such that 
\(P = P_{n-1} \cdot \dots \cdot P_1\) for which 

\[
    PA = LU
\]

where \(L\) and \(U\) are lower and upper triangular matrices respectively.

\subsubsection{Extended \(LU\)-Decomposition Algorithm}

The only important change for the algorithm is that for each row swaps to our matrix \(A^{(i)}\) we perform the same on the matrix 
\(P^{(i)}\). The process for the Gaussian elimination remains the same.

Overview:

\begin{enumerate}

    \item Decompose \(A\) into  \(L,U\) and \(P\).

    \item Solve \(Ly = Pb\). By first permuting \(b\) just like \(A\) and then \(Ly = Pv\)
    
    \item Solve \(Ux = y\). 

\end{enumerate}

\textbf{Example:}

Find the \(L,U\) and \(P\) for

\[
    A = 
    \begin{bmatrix}
        0 & 0 & 0 & 1 \\ 
        2 & 1 & 2 & 0 \\ 
        4 & 4 & 0 & 0 \\ 
        2 & 3 & 1 & 0
    \end{bmatrix}
\]

We now use a vector \(p\) of size \(n \times 1\) where the \(j\)-th entry \(p_j\) is read as: switch the rows \(j\) with \(p_j\)

We swap rows 1 and 2, giving us \(p_1 = 2\):

\[
    \begin{bmatrix} 
        2 & 1 & 2 & 0 \\ 
        0 & 0 & 0 & 1 \\ 
        4 & 4 & 0 & 0 \\ 
        2 & 3 & 1 & 0
    \end{bmatrix}
\]

Proceeding with the algorithm as for the simple \(LU\)-decomposition 

\[
    \begin{bmatrix} 
        2 & 1 & 2 & 0 \\ 
        0 & 0 & 0 & 1 \\ 
        4 & 4 & 0 & 0 \\ 
        2 & 3 & 1 & 0
    \end{bmatrix}
    \rightarrow
    \begin{bmatrix} 
        2 & 1 & 2 & 0 \\ 
        0 & 0 & 0 & 1 \\ 
        2 & 2 & -4 & 0 \\ 
        1 & 2 & -1 & 0
    \end{bmatrix}
    (p_2 = 3)
    \rightarrow
     \begin{bmatrix} 
        2 & 1 & 2 & 0 \\  
        2 & 2 & -4 & 0 \\ 
        0 & 0 & 0 & 1 \\ 
        1 & 2 & -1 & 0
    \end{bmatrix}
    \rightarrow
\]

\[
    \begin{bmatrix} 
        2 & 1 & 2 & 0 \\  
        2 & 2 & -4 & 0 \\ 
        0 & 0 & 0 & 1 \\ 
        1 & 1 & -1 & 0
    \end{bmatrix}
    \rightarrow
    \begin{bmatrix} 
        2 & 1 & 2 & 0 \\  
        2 & 2 & -4 & 0 \\ 
        0 & 0 & 0 & 1 \\ 
        1 & 1 & 3 & 0
    \end{bmatrix}
    (p_3 = 4)
    \rightarrow
    \begin{bmatrix} 
        2 & 1 & 2 & 0 \\  
        2 & 2 & -4 & 0 \\  
        1 & 1 & 3 & 0 \\
        0 & 0 & 0 & 1 \\ 
    \end{bmatrix}
\]

Hence 

\[
    L = 
    \begin{bmatrix} 
        1 & 0 & 0 & 0 \\  
        2 & 1 & 0 & 0 \\  
        1 & 1 & 1 & 0 \\
        0 & 0 & 0 & 1 \\ 
    \end{bmatrix}
    \quad 
    U = 
    \begin{bmatrix} 
        2 & 1 & 2 & 0 \\  
        0 & 2 & -4 & 0 \\  
        0 & 0 & 3 & 0 \\
        0 & 0 & 0 & 1 \\ 
    \end{bmatrix}
    \quad 
    v = 
    \begin{bmatrix}
       2 \\ 
       3 \\
       4 
    \end{bmatrix}
\]

Now the solving part is equal for both procedures: 

First \(Ly = Pb\), where \(Pb\) is the permuted vector using \(p\).

\[
    \begin{bmatrix} 
        1 & 0 & 0 & 0 \\  
        2 & 1 & 0 & 0 \\  
        1 & 1 & 1 & 0 \\
        0 & 0 & 0 & 1 \\ 
    \end{bmatrix}
    y = 
    \begin{bmatrix}
        5 \\ 
        4 \\ 
        5 \\ 
        3
    \end{bmatrix}
    \iff 
    y = 
    \begin{bmatrix}
        5 \\ 
        -6 \\
        6 \\ 
        3  
    \end{bmatrix}
\]

Finally, \(Ux = y\):

\[
    \begin{bmatrix} 
        2 & 1 & 2 & 0 \\  
        0 & 2 & -4 & 0 \\  
        0 & 0 & 3 & 0 \\
        0 & 0 & 0 & 1 \\ 
    \end{bmatrix}
    x = 
    \begin{bmatrix}
        5 \\ 
        -6 \\
        6 \\ 
        3  
    \end{bmatrix}
    \iff
    x =
    \begin{bmatrix}
        0 \\ 
        1 \\ 
        2 \\
        3
    \end{bmatrix}
\]

\subsubsection*{Greatest Pivot Strategy}

As a side note, to mitigate the rounding errors in scientific computing, the swap row strategy is 
modified to always take the row with biggest element in the current column.

\subsection{Cholesky Decomposition}

This is a special a special case of the \(LU\)-Decomposition. Given \(A \in \R^{n \times n}\) a symmetric, positive definite matrix it can be decomposed as 

\[
    A = \tilde{L}D\tilde{L}^T,
\]

where \(\tilde{L}\) is a lower triangular matrix with 1's in the main diagonal and \(D\) a diagonal matrix with positive 
diagonal elements. (\emph{Rational Cholesky Decomposition}).

Or \(A = G G^T\), \(G\) is lower triangular matrix. (\emph{Cholesky Decomposition}).

\subsubsection{Algorithm for the Cholesky Decomposition}

\begin{enumerate}

    \item Find the \(LU\)-decomposition of \(A\). 

    \item \(D\) is given by the entries in the diagonal of \(U\). 
    
    \item Determine \(V\) by taking \(L^T\), replacing its main diagonal with the diagonal of \(U\) and then dividing each row with the pivot 
    elements of \(D\). For symmetric matrices \(V = L^T\).

    \item Find the root of \(D\) by splitting it into two diagonal matrices which consist on the square roots of the elements in the main diagonal.
    Take one of them and call it \(S\).

    \item Find \(G = LS\). 
    
    \item Finally \(A = G G^T\).

\end{enumerate}

\textbf{Example:}

Find the Cholesky decomposition of 

\[
    A = 
    \begin{bmatrix}
       1 & 1 & 1 \\ 
       1 & 2 & -2 \\ 
       1 & -2 & 14 
    \end{bmatrix}
    \rightarrow 
    \begin{bmatrix}
       1 & 1 & 1 \\ 
       1 & 2 & -2 \\ 
       1 & -2 & 14 
    \end{bmatrix}
    \rightarrow 
    \begin{bmatrix}
       1 & 1 & 1 \\ 
       1 & 1 & -3 \\ 
       1 & -3 & 13 
    \end{bmatrix}
\]

\[
    \begin{bmatrix}
       1 & 1 & 1 \\ 
       1 & 1 & -3 \\ 
       1 & -3 & 13 
    \end{bmatrix}
    \rightarrow 
    \begin{bmatrix}
       1 & 1 & 1 \\ 
       1 & 1 & -3 \\ 
       1 & -3 & 4 
    \end{bmatrix}
\]

We get 

\[
    L = 
    \begin{bmatrix}
       1 & 0 & 0 \\ 
       1 & 1 & 0 \\ 
       1 & -3 & 1 
    \end{bmatrix}
    \quad 
    U = 
    \begin{bmatrix}
       1 & 1 & 1 \\ 
       0 & 1 & -3 \\ 
       0 & 0 & 4 
    \end{bmatrix}
    \quad 
    D = 
    \begin{bmatrix}
       1 & 0 & 0 \\ 
       0 & 1 & 0 \\ 
       0 & 0 & 4 
    \end{bmatrix}
\]

Then, we get 

\[
    S = 
    \begin{bmatrix}
       1 & 0 & 0 \\ 
       0 & 1 & 0 \\ 
       0 & 0 & 2 
    \end{bmatrix}
    \quad 
    V = 
    \begin{bmatrix}
       1 & 1 & 1 \\ 
       0 & 1 & -3 \\ 
       0 & 0 & 2 
    \end{bmatrix}
\]

And finally, 

\[
    Q = LS = 
     \begin{bmatrix}
       1 & 0 & 0 \\ 
       1 & 1 & 0 \\ 
       1 & -3 & 1 
    \end{bmatrix} 
    \begin{bmatrix}
       1 & 0 & 0 \\ 
       0 & 1 & 0 \\ 
       0 & 0 & 2 
    \end{bmatrix}
    = 
    \begin{bmatrix}
       1 & 0 & 0 \\ 
       1 & 1 & 0 \\ 
       1 & -3 & 2 
    \end{bmatrix}
\]

We got \(A = Q Q^T\).
